%
% Copyright (c) 2015 Radoslaw Kujawa.
% All rights reserved.
%
% Redistribution and use in source and binary forms, with or without
% modification, are permitted provided that the following conditions
% are met:
%
% 1. Redistributions of source code must retain the above copyright
%    notice, this list of conditions and the following disclaimer.
% 2. Redistributions in binary form must reproduce the above copyright
%    notice, this list of conditions and the following disclaimer in the
%    documentation and/or other materials provided with the distribution.
%
% THIS SOFTWARE IS PROVIDED BY RADOSLAW KUJAWA (THE AUTHOR) AND CONTRIBUTORS
% ``AS IS'' AND ANY EXPRESS OR IMPLIED WARRANTIES, INCLUDING, BUT NOT LIMITED
% TO, THE IMPLIED WARRANTIES OF MERCHANTABILITY AND FITNESS FOR A PARTICULAR
% PURPOSE ARE DISCLAIMED.  IN NO EVENT SHALL THE AUTHOR OR CONTRIBUTORS
% BE LIABLE FOR ANY DIRECT, INDIRECT, INCIDENTAL, SPECIAL, EXEMPLARY, OR
% CONSEQUENTIAL DAMAGES (INCLUDING, BUT NOT LIMITED TO, PROCUREMENT OF
% SUBSTITUTE GOODS OR SERVICES; LOSS OF USE, DATA, OR PROFITS; OR BUSINESS
% INTERRUPTION) HOWEVER CAUSED AND ON ANY THEORY OF LIABILITY, WHETHER IN
% CONTRACT, STRICT LIABILITY, OR TORT (INCLUDING NEGLIGENCE OR OTHERWISE)
% ARISING IN ANY WAY OUT OF THE USE OF THIS SOFTWARE, EVEN IF ADVISED OF THE
% POSSIBILITY OF SUCH DAMAGE.
%
% 
\documentclass[dvipsnames,table]{beamer}
\usepackage{polski}

\usetheme{Rochester}
\usecolortheme{orchid}

\usepackage{listings}
\usepackage{ucs}
\usepackage[utf8x]{inputenc}
\usepackage{wasysym}
\usepackage[normalem]{ulem}
\usepackage{amsmath}
\usepackage{hyperref}
\usepackage{tikzsymbols}

\setbeamertemplate{navigation symbols}{}
\setbeamertemplate{caption}[numbered]
\setbeamerfont{caption}{size=\scriptsize}
\setbeamercolor{framenote}{bg=OSEC-red!25}
\setbeamercolor{rednote}{bg=Red!25}
\setbeamercolor{palette primary}{use=structure,fg=white,bg=OSEC-red}
\setbeamercolor{palette secondary}{use=structure,fg=white,bg=OSEC-red2}

\setbeamertemplate{itemize item}{\scriptsize\raise1pt\hbox{\donotcoloroutermaths$\blacktriangleright$}}
\setbeamertemplate{itemize subitem}{\tiny\raise1pt\hbox{\donotcoloroutermaths$\bullet$}}
\setbeamertemplate{itemize subsubitem}{\tiny\raise1pt\hbox{\donotcoloroutermaths{--}}}

\setbeamertemplate{enumerate item}{\insertenumlabel.}
\setbeamertemplate{enumerate subitem}{\insertenumlabel.\insertsubenumlabel}
\setbeamertemplate{enumerate subsubitem}{\insertenumlabel.\insertsubenumlabel.\insertsubsubenumlabel}
\setbeamertemplate{enumerate mini template}{\insertenumlabel}

\setbeamercolor{itemize item}{fg=OSEC-red, bg=OSEC-red}
\setbeamercolor{itemize subitem}{fg=OSEC-red, bg=OSEC-red}
\setbeamercolor{itemize subsubitem}{fg=OSEC-red, bg=OSEC-red}

\setbeamercolor{section number projected}{fg=white,bg=OSEC-red}
\setbeamercolor{subsection number projected}{fg=white,bg=OSEC-red}
\setbeamercolor{button}{bg=OSEC-red,fg=white}

\setbeamertemplate{section in toc}[circle]
\setbeamertemplate{subsection in toc}[square]


\definecolor{OSEC-red}{RGB}{160,29,44}
\definecolor{OSEC-red2}{RGB}{177,76,12}
\hypersetup{colorlinks=true,linkcolor=white,urlcolor=OSEC-red}

\setlength{\tabcolsep}{8pt}
\renewcommand{\arraystretch}{1.2}

\newcommand{\nbsdcolor}[1] {
	{\color{OSEC-red} #1}
}

\lstset{
   language=java,
   basicstyle=\tiny,
   breaklines=true,
   escapechar=\@,
   commentstyle=\color{OSEC-red}
}

\AtBeginSection[]{
\frame{

\begin{center}

{\usebeamerfont{section name}\usebeamercolor[fg]{section name}Część~\insertsectionnumber}
    \vskip1em\par

	\begin{beamercolorbox}[sep=12pt,center]{palette primary}
		\usebeamerfont{section title}\insertsection\par
	\end{beamercolorbox}
\end{center}

%\sectionpage
}
}

\title{Wykorzystanie rekordów DNS \\ do weryfikacji kluczy publicznych SSH}


\author{Radosław Kujawa -- radoslaw.kujawa@osec.pl}

\institute{OSEC}

\begin{document}

\begin{frame}
\titlepage
\end{frame}

\begin{frame}[allowframebreaks]
\frametitle{Spis treści}
{
\hypersetup{colorlinks=true,linkcolor=black,urlcolor=OSEC-red}
\tableofcontents
}
\end{frame}

\section{Zagrożenia bezpieczeństwa związane z weryfikacją kluczy publicznych hostów}

\begin{frame}
\frametitle{SSH - weryfikacja kluczy a bezpieczeństwo}
\begin{itemize}
	\item Każdy z nas używa SSH...
	\item Protokół SSH jest ,,bezpieczny''.
	\item Ale często zapominamy o zagrożeniach.
\end{itemize}
\end{frame}

\begin{frame}[fragile]
\frametitle{Akceptacja klucza publicznego zdalnego hosta}
\scriptsize
\begin{verbatim}
$ ssh example.com
The authenticity of host 'example.com (1.2.3.4)' can't be established.
RSA key fingerprint is 12:75:dd:41:22:3f:6f:88:f7:bc:4e:ad:2e:b5:da:2d.
Are you sure you want to continue connecting (yes/no)?
\end{verbatim}
\normalsize
\begin{itemize}
	\item Manualna weryfikacja odcisku klucza najczęściej się nie odbywa.
	\begin{itemize}
		\item A co to jest? Myślałem, że trzeba tu wpisać ,,yes''.
		\item Nie mam czasu / zarobiony jestem.
		\item A kto będzie atakował ten mało ważny serwer?
	\end{itemize}

\end{itemize}

\end{frame}

\begin{frame}
\frametitle{Problem z weryfikacją kluczy publicznych}
\begin{itemize}
	\item Protokół SSH jest faktycznie bezpieczny, ale...
	\item Całe bezpieczeństwo opiera się na weryfikacji klucza publicznego zdalnego hosta!
	\item Skutecznie wykonany atak Man-in-the-Middle podczas pierwszego połączenia całkowicie burzy to bezpieczeństwo.
	\item SSH nie wymusza stosowania infrastruktury klucza publicznego (PKI).
\end{itemize}
\end{frame}

\begin{frame}
\frametitle{Gdy klucz zostanie zaakceptowany}
\begin{itemize}
	\item Jest zapisywany w {\tt \$HOME/.ssh/known\_hosts}.
	\item Używany do weryfikacji każdego kolejnego połączenia do tego samego hosta.
	\item Przy połączeniu do nowego hosta problem powraca...
\end{itemize}
\end{frame}

\section{Jak wykorzystać rekordy SSHFP do zautomatyzowanej weryfikacji kluczy}

\begin{frame}
\frametitle{Mechanizm SSHFP}
\begin{itemize}
	\item \href{https://tools.ietf.org/html/rfc4255}{RFC4255},
		  \href{https://tools.ietf.org/html/rfc6594}{RFC6594},
		  \href{https://tools.ietf.org/html/rfc7479}{RFC7479}.
	\item Klient inicjuje połączenie z serwerem SSH.
	\item Serwer SSH legitymuje się swoim kluczem publicznym.
	\item Klient pobiera wartość rekordów SSHFP dla danego FQDN.
	\item Klient weryfikuje odcisk klucza dla danego algorytmu, skoro odcik się zgadza...
	\begin{itemize}
		\item Może automatycznie zaakceptować klucz.
		\item Może zapytać użytkownika czy klucz zaakceptować, podając informację o zgodności z rekordem SSHFP.
	\end{itemize}
	\item Dalej połączenie SSH przebiega jak zazwyczaj.
\end{itemize}
\end{frame}

\begin{frame}
\frametitle{Generowanie rekordów SSHFP}
\begin{itemize}
	\item Nie trzeba niczego instalować dodatkowo \Smiley
	\item Znane narzędzie {\tt ssh-keygen} posiada tą funkcjonalność.
	\item {\tt ssh-keygen -r `hostname`}.
\end{itemize}
\end{frame}

\begin{frame}[fragile]
\frametitle{Struktura rekordu SSHFP}
\scriptsize
\begin{verbatim}
www		IN A	 192.168.0.100
www		IN SSHFP 1 1 0b34980da528297a8c9f3be3e8b6c05d58d14602
www		IN SSHFP 1 2 1f45321c7f1b730840a7fef14e30725a9d5ef5c5c7a9174(...)
www		IN SSHFP 2 1 c9a822aba840658002a4911f0a430f00f6e76d4e
www		IN SSHFP 2 2 15d544b67d0074bcd74219be55b96cdb36378f513159582(...)
\end{verbatim}
\normalsize
\begin{itemize}
	\item Algorytm którym wygenerowano klucz:
	\begin{itemize}
		\item {\tt 1} - RSA.
		\item {\tt 2} - DSA.
		\item {\tt 3} - ECDSA.
		\item {\tt 4} - Ed25519.
	\end{itemize}
	\item Funkcja skrótu którą wygenerowano odcisk:
	\begin{itemize}
		\item {\tt 1} - SHA1.
		\item {\tt 2} - SHA256.
	\end{itemize}
	\item Weryfikacja: {\tt dig -t SSHFP www.example.com}.
\end{itemize}
\end{frame}

\section{Konfiguracja klienta SSH}

\begin{frame}
\frametitle{Konfiguracja klienta -- OpenSSH}
\begin{itemize}
	\item Nie trzeba niczego instalować dodatkowo \Smiley.
	\item Parametry konfiguracyjne klienta:
	\begin{itemize}
		\item {\tt /etc/ssh/ssh\_config}.
		\item {\tt \$HOME/.ssh/config}.
		\item Parametr {\tt -o} polecenia {\tt ssh}. 
	\end{itemize}
	\item Opcje:
	\begin{itemize}
		\item {\tt VerifyHostKeyDNS} na wartość {\tt yes} lub {\tt ask}.
	\end{itemize}
\end{itemize}
\end{frame}

\section{Automatyzacja procesu aktualizacji stref}

\begin{frame}
\frametitle{Automatyzacja aktualizacji stref}
\begin{itemize}
	\item Brak standardu \Sadey.
	\item Red Hat IPA/FreeIPA.
	\item Dynamic DNS -- {\tt nsupdate}.
	\item Generowanie kluczy w Kickstarcie i własny skrypt do aktualizacji DNS.
	\item Puppet/Chef/Ansible...
\end{itemize}
\end{frame}

\section{Zabezpieczanie DNS przed spoofingiem}

\begin{frame}
\frametitle{DNS a spoofing}
\begin{itemize}
	\item Brak podpisu elektronicznego DNS w standardzie \Sadey.
	\begin{itemize}
		\item Podatność DNS na ataki MITM.
	\end{itemize}
	\item DNSSEC rozwiązuje ten problem.
	\begin{itemize}
		\item Tylko trzeba go wdrożyć...
	\end{itemize}
	\item W obrębie własnej infrastruktury alternatywą może być IPsec.
	
	\item Nawet jeśli ,,spoofowalność'' DNS jest problemem...
	\begin{itemize}
		\item Samo wdrożenie SSHFP jest lepsze w porównaniu z brakiem jakiejkolwiek weryfikacji kluczy publicznych hostów.
	\end{itemize}
\end{itemize}
\end{frame}

\begin{frame}
\frametitle{Alternatywy dla SSHFP?}
\begin{itemize}
	\item OpenSSH CA (min. wersja 5.3, faktycznie potrzebna 5.6).
	\item {\tt sss\_ssh\_knownhostsproxy}.
	\item {\tt ssh\_keyscan}.
	\item Manualna weryfikacja \Sadey.
\end{itemize}
\end{frame}

\begin{frame}
\frametitle{Koniec\ldots}
\begin{center}
\includegraphics[scale=0.5]{img-oseclogo.png}

Dziękuje!

Czy są pytania?

\end{center}

\end{frame}

 
\end{document}
